## 📚 ПОЛНЫЙ СПИСОК ДОМЕНОВ 

1. ПСИХОЛОГИЯ ?
2. НЕЙРОНАУКИ ?
3. СОЦИОЛОГИЯ 
4. ЭКОНОМИКА 
5. БИЗНЕС И УПРАВЛЕНИЕ 
6. ФИЛОСОФИЯ 
7. ЛИНГВИСТИКА
8. ИСКУССТВО (art)
9. МУЗЫКА
10. АНТРОПОЛОГИЯ 
11. РЕЛИГИОВЕДЕНИЕ
12. МИФОЛОГИЯ
13. ИСТОРИЯ 
14. ПЕДАГОГИКА 
15. КОММУНИКАЦИЯ 
16. СЕМИОТИКА 
17. МЕДИАЛОГИЯ
18. РИТОРИКА
19. ЮРИСПРУДЕНЦИЯ
20. ПОЛИТОЛОГИЯ 
21. ДИПЛОМАТИЯ 
22. КОНФЛИКТОЛОГИЯ 
23. ДИЗАЙН
24. ХОРЕОГРАФИЯ
25. ДРАМАТУРГИЯ 
26. КИНЕМАТОГРАФИЯ
27. ГЕЙМ-ДИЗАЙН 
28. СПОРТИВНАЯ НАУКА
29. СОМАТИКА 
30. ТАНЦЕВАЛЬНО-ДВИГАТЕЛЬНАЯ ТЕРАПИЯ
31. ЭРГОНОМИКА
32. ПСИХОТЕРАПИЯ ?
33. КОУЧИНГ 
34. МЕДИТАЦИЯ И MINDFULNESS 
35. ХОЛИСТИЧЕСКИЕ ПРАКТИКИ
36. ЭКОЛОГИЯ ЧЕЛОВЕКА
37. УРБАНИСТИКА
38. ПЕРМАКУЛЬТУРА
39. СИСТЕМНАЯ БИОЛОГИЯ
40. БИБЛИОТЕКОВЕДЕНИЕ
41. АРХИВОВЕДЕНИЕ
42. DATA SCIENCE 
43. UX/UI ДИЗАЙН 
44. КУЛИНАРИЯ
45. РИТУАЛИСТИКА 
46. ОНЕЙРОЛОГИЯ (НАУКА О СНАХ)
47. ПАРАДОКСОЛОГИЯ 
48. СИСТЕМНОЕ МЫШЛЕНИЕ 

## 📚 ПОЛНЫЙ СПИСОК ДОМЕНОВ И ПОДДОМЕНОВ ДЛЯ АНАЛИЗА ФУНКТОРОВ

### 1. ПСИХОЛОГИЯ

- Когнитивная психология
- Социальная психология
- Психология развития
- Клиническая психология
- Психология личности
- Нейропсихология
- Психология восприятия
- Психология памяти
- Психология эмоций
- Психология мотивации
- Психология сознания
- Трансперсональная психология

### 2. НЕЙРОНАУКИ

- Молекулярная нейробиология
- Системная нейронаука
- Когнитивная нейронаука
- Вычислительная нейронаука
- Нейропластичность
- Нейрофармакология
- Нейроэндокринология
- Нейроиммунология

### 3. СОЦИОЛОГИЯ

- Микросоциология
- Макросоциология
- Социология культуры
- Социология организаций
- Социология девиантного поведения
- Социология социальных движений
- Социология семьи
- Цифровая социология

### 4. ЭКОНОМИКА

- Микроэкономика
- Макроэкономика
- Поведенческая экономика
- Экономика развития
- Институциональная экономика
- Экспериментальная экономика
- Нейроэкономика
- Эконофизика

### 5. БИЗНЕС И УПРАВЛЕНИЕ

- Стратегический менеджмент
- Операционный менеджмент
- Управление изменениями
- Организационное поведение
- Маркетинг
- Управление инновациями
- Предпринимательство
- Управление проектами

### 6. ФИЛОСОФИЯ

- Эпистемология
- Онтология
- Этика
- Эстетика
- Философия сознания
- Философия языка
- Философия науки
- Феноменология
- Экзистенциализм
- Прагматизм

### 7. ЛИНГВИСТИКА

- Фонетика и фонология
- Морфология
- Синтаксис
- Семантика
- Прагматика
- Социолингвистика
- Психолингвистика
- Нейролингвистика
- Компьютерная лингвистика

### 8. ИСКУССТВО

- Живопись
- Скульптура
- Инсталляция
- Перформанс
- Цифровое искусство
- Концептуальное искусство
- Ленд-арт
- Биоарт

### 9. МУЗЫКА

- Теория музыки
- Композиция
- Импровизация
- Музыкальная акустика
- Психоакустика
- Этномузыкология
- Электронная музыка
- Алгоритмическая композиция

### 10. АНТРОПОЛОГИЯ

- Культурная антропология
- Социальная антропология
- Физическая антропология
- Лингвистическая антропология
- Визуальная антропология
- Киберантропология
- Медицинская антропология
- Экологическая антропология

### 11. РЕЛИГИОВЕДЕНИЕ

- Сравнительное религиоведение
- Феноменология религии
- Психология религии
- Социология религии
- Мистицизм
- Ритуаловедение
- Священные тексты
- Новые религиозные движения

### 12. МИФОЛОГИЯ

- Сравнительная мифология
- Структурная мифология
- Архетипы
- Героический эпос
- Космогонические мифы
- Эсхатология
- Фольклор
- Городские легенды

### 13. ИСТОРИЯ

- История цивилизаций
- История технологий
- История идей
- Военная история
- История науки
- История искусства
- Микроистория
- История повседневности

### 14. ПЕДАГОГИКА

- Дидактика
- Теория обучения
- Андрагогика
- Специальная педагогика
- Цифровая педагогика
- Критическая педагогика
- Монтессори-педагогика
- Вальдорфская педагогика

### 15. КОММУНИКАЦИЯ

- Теория коммуникации
- Межличностная коммуникация
- Массовая коммуникация
- Невербальная коммуникация
- Межкультурная коммуникация
- Организационная коммуникация
- Кризисная коммуникация
- Цифровая коммуникация

### 16. СЕМИОТИКА

- Структурная семиотика
- Когнитивная семиотика
- Социосемиотика
- Биосемиотика
- Киберсемиотика
- Визуальная семиотика
- Семиотика культуры
- Семиотика медиа

### 17. МЕДИАЛОГИЯ

- Теория медиа
- Трансмедиа
- Новые медиа
- Медиаэкология
- Медиаархеология
- Медиафилософия
- Визуальные исследования
- Sound studies

### 18. РИТОРИКА

- Классическая риторика
- Новая риторика
- Визуальная риторика
- Цифровая риторика
- Политическая риторика
- Судебная риторика
- Академическая риторика
- Корпоративная риторика

### 19. ЮРИСПРУДЕНЦИЯ

- Теория права
- Конституционное право
- Международное право
- Правовая антропология
- Социология права
- Психология права
- Кибернетика права
- Правовая герменевтика

### 20. ПОЛИТОЛОГИЯ

- Политическая теория
- Сравнительная политология
- Политическая психология
- Политическая социология
- Политическая экономия
- Геополитика
- Электоральные исследования
- Политическая коммуникация

### 21. ДИПЛОМАТИЯ

- Классическая дипломатия
- Публичная дипломатия
- Цифровая дипломатия
- Превентивная дипломатия
- Экономическая дипломатия
- Культурная дипломатия
- Парадипломатия
- Медиация

### 22. КОНФЛИКТОЛОГИЯ

- Теория конфликта
- Управление конфликтами
- Разрешение конфликтов
- Медиация
- Переговоры
- Миротворчество
- Постконфликтное восстановление
- Превенция конфликтов

### 23. ДИЗАЙН

- Графический дизайн
- Промышленный дизайн
- Дизайн интерьера
- Ландшафтный дизайн
- Веб-дизайн
- UX/UI дизайн
- Сервис-дизайн
- Спекулятивный дизайн

### 24. ХОРЕОГРАФИЯ

- Классический танец
- Современный танец
- Народный танец
- Контактная импровизация
- Буто
- Соматические практики
- Танцевальная композиция
- Лабан-анализ движения

### 25. ДРАМАТУРГИЯ

- Классическая драматургия
- Современная драматургия
- Сценография
- Режиссура
- Актёрское мастерство
- Театральная педагогика
- Перформативные исследования
- Постдраматический театр

### 26. КИНЕМАТОГРАФИЯ

- Теория кино
- Монтаж
- Кинодраматургия
- Документальное кино
- Экспериментальное кино
- Анимация
- Видеоарт
- VR-кинематография

### 27. ГЕЙМ-ДИЗАЙН

- Игровая механика
- Нарративный дизайн
- Level design
- Балансировка
- Монетизация
- Геймификация
- Serious games
- AR/VR игры

### 28. СПОРТИВНАЯ НАУКА

- Биомеханика
- Физиология спорта
- Психология спорта
- Спортивная медицина
- Теория тренировки
- Тактика и стратегия
- Спортивная педагогика
- Адаптивный спорт

### 29. СОМАТИКА

- Метод Фельденкрайза
- Техника Александера
- Бодинамика
- Розен-метод
- Continuum Movement
- Body-Mind Centering
- Аутентичное движение
- Идеокинезис

### 30. ТАНЦЕВАЛЬНО-ДВИГАТЕЛЬНАЯ ТЕРАПИЯ

- Психодинамический подход
- Гуманистический подход
- Медицинская модель
- Интегративный подход
- Развивающий подход
- Травма-информированный подход
- Групповая ТДТ
- Семейная ТДТ

### 31. ЭРГОНОМИКА

- Физическая эргономика
- Когнитивная эргономика
- Организационная эргономика
- Макроэргономика
- Нейроэргономика
- Эргономика интерфейсов
- Эргономика рабочего места
- Эргономика транспорта

### 32. ПСИХОТЕРАПИЯ

- Психоанализ
- Гештальт-терапия
- Когнитивно-поведенческая терапия
- Экзистенциальная терапия
- Системная семейная терапия
- Арт-терапия
- EMDR
- Соматическая терапия
- Процессуальная терапия
- Нарративная терапия
- ACT
- DBT

### 33. КОУЧИНГ

- Лайф-коучинг
- Бизнес-коучинг
- Карьерный коучинг
- Трансформационный коучинг
- Интегральный коучинг
- Онтологический коучинг
- Соматический коучинг
- Нейрокоучинг

### 34. МЕДИТАЦИЯ И MINDFULNESS

- Випассана
- Дзен
- Трансцендентальная медитация
- MBSR
- MBCT
- Метта-медитация
- Медитация движения
- Нейрофидбек

### 35. ХОЛИСТИЧЕСКИЕ ПРАКТИКИ

- Аюрведа
- Традиционная китайская медицина
- Гомеопатия
- Натуропатия
- Остеопатия
- Хиропрактика
- Энергетическая медицина
- Интегративная медицина

### 36. ЭКОЛОГИЯ ЧЕЛОВЕКА

- Экологическая психология
- Биофилия
- Экотерапия
- Глубинная экология
- Социальная экология
- Экология здоровья
- Экология питания
- Экология жилища

### 37. УРБАНИСТИКА

- Городское планирование
- Транспортное планирование
- Социология города
- Психогеография
- Тактический урбанизм
- Умные города
- Партиципаторное проектирование
- Городская антропология

### 38. ПЕРМАКУЛЬТУРА

- Дизайн систем
- Зоны и сектора
- Пищевые леса
- Водные системы
- Почвоведение
- Социальная пермакультура
- Экономика пермакультуры
- Городская пермакультура

### 39. СИСТЕМНАЯ БИОЛОГИЯ

- Метаболомика
- Протеомика
- Геномика
- Эпигенетика
- Системная медицина
- Синтетическая биология
- Биоинформатика
- Сетевая биология

### 40. БИБЛИОТЕКОВЕДЕНИЕ

- Организация знаний
- Информационный поиск
- Метаданные
- Цифровые библиотеки
- Библиометрия
- Управление знаниями
- Информационная грамотность
- Сохранение цифрового наследия

### 41. АРХИВОВЕДЕНИЕ

- Теория архивов
- Архивная эвристика
- Электронные архивы
- Оценка документов
- Консервация
- Архивная педагогика
- Общественные архивы
- Память и идентичность

### 42. DATA SCIENCE

- Машинное обучение
- Глубокое обучение
- Обработка естественного языка
- Компьютерное зрение
- Анализ временных рядов
- Графовая аналитика
- Причинный вывод
- Объяснимый ИИ

### 43. UX/UI ДИЗАЙН

- Исследование пользователей
- Информационная архитектура
- Проектирование взаимодействия
- Визуальный дизайн интерфейсов
- Прототипирование
- Юзабилити-тестирование
- Доступность
- Дизайн-системы

### 44. КУЛИНАРИЯ

- Молекулярная гастрономия
- Ферментация
- Выпечка
- Соусы и эмульсии
- Термическая обработка
- Консервирование
- Food pairing
- Кулинарная антропология

### 45. РИТУАЛИСТИКА

- Обряды перехода
- Календарные ритуалы
- Целительские ритуалы
- Погребальные обряды
- Инициации
- Ритуалы власти
- Современные ритуалы
- Ритуальное пространство

### 46. ОНЕЙРОЛОГИЯ (НАУКА О СНАХ)

- Физиология сна
- Психология сновидений
- Осознанные сновидения
- Сновидческие практики культур
- Толкование снов
- Кошмары и их терапия
- Гипнагогия
- Инкубация снов

### 47. ПАРАДОКСОЛОГИЯ

- Логические парадоксы
- Семантические парадоксы
- Парадоксы бесконечности
- Квантовые парадоксы
- Парадоксы времени
- Парадоксы выбора
- Парадоксы восприятия
- Социальные парадоксы

### 48. СИСТЕМНОЕ МЫШЛЕНИЕ

- Общая теория систем
- Системная инженерия
- Системная динамика
- Кибернетика
- Сложные адаптивные системы 
- Методология мягких систем (Soft Systems)
- Системная экология
- Системные архетипы
- Ментальные модели
- Петли обратной связи
- Устойчивость и устойчивые системы